\section{The Canonical Natural-Number System}
\label{sec:canonical-natural-number-system}

\begin{topicbox}{Working in \texorpdfstring{$\mathbb{N}$}{N}}
\begin{exposition}
The preceding chapter established the Peano-system facts needed for arithmetic:
induction, successor analysis, recursion, and categoricity. From this point
forward, $\mathbb{N}$ denotes the canonical natural-number Peano system, with
distinguished first element $1$ and successor map $S$. The arithmetic
operations introduced in this chapter are defined inside that system by the
recursion principles already proved for Peano systems.
\end{exposition}

\begin{remark*}[Standing convention]
Unless otherwise stated, variables $m,n,k$ range over $\mathbb{N}$, the symbol
$1$ denotes the distinguished first natural number, and $S(n)$ denotes the
successor of $n$.
\end{remark*}
\end{topicbox}
